\documentclass{article}

\usepackage[utf8]{inputenc}

\title{Title}
\author{Author}
\date{\today}

\begin{document}

\maketitle

\section{Heilmeier}

\begin{enumerate}
    \item What are you trying to do?
    \begin{enumerate}
        \item Provide accurate and feasible local, context-sensitive causal explanations in a causal probabilistic machine learning model, and, by the same token, provide explanations for events in the world, insofar as the model adequately represents the relevant factors.
        
        \item Causal responsibility attribution is important: for example, we might want to understand why an adverse outcome occurred, in order to prevent it from happening again later; or we might want to understand exactly how and why a predicted, counterfactual positive outcome can be brought about.
    \end{enumerate}
    \item Why is it hard right now?
    \begin{enumerate}
        \item Most of the existing scalable explanatory methods that could be applied to a model are not causal.
        \item Current methods for causal responsibility attribution often involve combinatorics and existential quantifiers that quickly become intractable, and so don't scale.
        \item Many of the existing definitions of causal responsibility and methods assume constrained parametric causal models, often assuming the relevant variables are binary or at least discrete. We need a tool that we can work independently on the type and support of the random variables at play. 
    \end{enumerate}
    \item What's new in your approach and why will it work?
    \begin{enumerate}
    \item Inspired by (i) the use of witness nodes in Halpern's AC, and (ii) the parallel world treatment present in Pearl's notion of PNS. The former allows us for context sensitivity and patterns of reasoning akin to generalized mediation analysis, and the latter
    \item AC: precedent for witness sets, and the modified definition for populating their values with the factual ones,  
    \item PNS for integration of exogenous noise, both use alternative values, but only for categorical variables, we allow for analogous treatment of continuous variables
    \item Normality gives precedent for excised distributions and potentially thinking in terms of distributions over contexts 
    \item sampling over suspects and witnesses
    \item excised distribution for alternative values (a probabilistic take on an explication of normality)
    \item sampling over subsets and conditioning on antecedent activation provides an alternative mechanism akin to minimality constraint in the ac def
    \item push-forward causal impact score for variables of different types and supports, including allowing for continuous causes, witnesses, and outcomes
    \item using integration instead of existential quantifier, with the idea that this is closer to normality considerations
    \item This is achieved by causal probabilistic model transformations, which allow us to turn search into an inference problem (refer to search as inference stuff).
     \item This allows for the use of well-established approximation algorithms, which makes approximate search for explanations computationally feasible even for large-scale causal probabilistic models (refer to the HouseIQ example)
        

        % \item We provide a framework for causal explanation based on separate notions of Sufficiency and Necessity by merging Pearl's Probability of Necessity and Sufficiency with Halpern's Actual Causality based on witnesses or context sets.
    \end{enumerate}
    \item If it works, who cares? What difference does it make?
    \begin{enumerate}
    \item The tooling has been implemented to allow use with large causal probabilistic models
    specified in ChiRho, 
    \item The formalism puts no strong requirements on the form of the model, form of the distributions, or form of the structural equations
    \item In result, it works for causal models whose components are vanilla ML tools, such as neural networks, or Gaussian processes
    \item Responsibility attribution has not yet been effectively mechanized in software for large, complicated models, which makes using exact notions present in the literature unfeasible. Our approach develops approximate causal attribution workable for users of models of non-trivial size.
    \end{enumerate}

    
    \item How do you know your thing works/solves the problem?
    
    \begin{enumerate}
        \item We have a theorem that shows that when used in specific settings with a restricted class of models, it recovers actual causality, with the only difference being that minimality is understood in terms of cardinality rather than subset relation (see theorems)
        \item We have benchmarked it against exact AC computation with a problem of scalable complexity, and while exact computations cease being feasible at models containing around 20 variables, the approximate methods can be used with much larger models, with accuracy constrained by MC sample sizes or whatever constraints come with the inference algorithm one uses with the transformed model
        \item We have successfully deployed the tool with a real-life spatio-temporal model involving neural network parametrized distributions (including normalizing flows) as causal kernels and a Gaussian Process component, trained on millions of datapoints.
        \item In that real-life setting, we benchmarked responsibility attribution against SHAP, which clearly does not give correct answers, in light of domain experts' opinions.
        \item We benchmarked the method using a synthetic run involving continuous variables in the context of causal overdetermination and undercutting
        \item We benchmarked the method in the context of a synthetic epidemiological problem involving overdetermination, undercutting, and a dynamical system (SIR) for disease modeling. Unlike simpler counterfactual queries, it correctly identifies causes of overshoot.
        \item We compared the behavior of our tooling to that of Differential Causal Effect in the context of a synthetic differentiable model, arguing that the sensitivity analyses that DCE enables are not always what we are after when we ask for local explanations, especially as DCE is not causally aware.

        \item Overall, people have been using non-causal explanation methods because they scaled and the causal ones didn't, until now.
        
        %the responsibility attributions given by the former are more causallstable, more intuitive, and more interpretable.
    \end{enumerate}
\end{enumerate}




\section{Relevant literature (topics)}


\begin{enumerate}

\item actual causality


\item probability of sufficiency and necessity

\item search as inference: \url{https://arxiv.org/abs/2406.17863}
% https://www.sciencedirect.com/science/article/abs/pii/S1364661312001957
% https://arxiv.org/abs/2406.17863
% https://arxiv.org/abs/2003.13221


\end{enumerate}

\section{Abstract}

\section{Intro}

\subsection{Paragraph Outline}

\begin{enumerate}
    \item Causal attribution/explanation is a fundamental reasoning task underpinning many important problems, such as establishing efficacy of medicines, pinpointing the drivers of economic growth, determining why a machine failed, and assessing the effects of public policy. (cite a bunch of papers using causal attribution generally, e.g. in mechanistic models used for policy, in deep machine learning models, in ... etc. a bunch of applications to demonstrate importance. include some non AI/ML examples where practitioners/policy makers/domain people are doing this ``in spirit'', but without formal tools.) Further, in the absence of scalable, formal causal attribution tools, people use ad-hoc and heuristic methods, like Granger causality, mechanistic explanations, correlations, or difference-in-differences. (citing the ad-hoc makes two points: (1) everyone wants causal attribution reasoning capabilities and (2) they aren't doing it well) (mention Granger causality). Even researchers with an advanced command of statistical methods may not possess a thorough understanding of how to assign causality, but would like such a tool.
We present a tool for assigning gradients of causality to a feature set, inspired by Pearl's Probability of Necessity and Sufficiency (PNS) based on counterfactuals and witness (context) sets. Our tool further extends PNS to provide local (variable-specific) explanations and gradients of attribution beyond binary causality prevalent in the literature.
Many researchers prefer a mechanistic explanation for cause, though this may encourage circular logic- using existing beliefs to find cause. Our software provides a direct tool which allows users to plug in data and immediately receive a score which provides the amount of causality for a given variable. 
The role of essentially all clinical drug trials is to establish whether a drug causes a cure.
*Drew's note: I like this excerpt "Arguably, in most statistical association tests, we, as researchers and clinicians, are ultimately interested in making causal inference. While statistical tests allow us to describe associations between two or more variables, the attempts to control for confounding and to adjust for the effect of other variables are all part of an agenda to infer causal effects."
https://onlinelibrary.wiley.com/doi/pdfdirect/10.1111/dmcn.14813
Practitioners across several fields have themselves noted the need for better frameworks to establish causality https://scite.ai/reports/causal-inference-and-observational-research-rODXOw
Causal Analysis on engineering using physical/mechanistic framework for identifying what caused an accident https://pdfs.semanticscholar.org/8a0a/714936bd21c52fd7f215c174f05eb7e25529.pdf
    \item ``Existing methods (used in stuff we cite above) either don't scale (AC, cite halpern's complexity analysis of AC) or, if they do scale, attribute impact to non-causal relationships, yielding misleading attributions and explanations (SHAP, etc.).''
Weak correlations can still be causal (Vitamin C and scurvy prevention, for example) and strong correlations can still be non-causal (like ice cream sales and forest fires).
Paper showing flaws in difference-in-differences https://papers.ssrn.com/sol3/papers.cfm?abstract_id=3315280
Paper showing flaws in Instrumental Variable https://ideas.repec.org/p/nbr/nberwo/5835.html#:~:text=association%20between%20quarter%20of%20birth,effect%20of%20education%20on%20earnings
Controlling for confounders (by adding them to your features) in a regression model tends to cause multicollinearity.
Mechanistic explanations can tend to support existing beliefs to validate claims- circular logic.
    \item Explain what our model does and give an overview of how it does it. Allude to later discussion on related work (``we take inspiration from'' ...), briefly comparing/contrasting existing approaches with what we do. Mention fenton-glynn/fang's stuff on probabilistic AC, mention responsibility and blame wrt taking expectations over exogenous noise distributions, mention normality as sometimes using probability to favor normal and over abnormal causes, pearl's probabilities of causation, etc.
    \item How do we test this/ensure that it makes sense? What experiments and case studies do we use.
    \item clearly enumerate claims/contributions
\end{enumerate}




 

\end{document}
